% Chapter 5: Chapter 5. Discussion


\section{FPGA Bring-up Challenges and Resolution}\label{sec:5_1}

The FPGA bring-up process for the E203-based CNN accelerator system revealed several non-trivial engineering challenges that are not typically encountered in pure RTL simulation environments. These challenges and their resolutions provide valuable insights for similar SoC prototyping efforts.


\subsection{Simulation-to-Hardware Gap}

A significant finding was that the RTL simulation, once properly configured, exactly reproduced the FPGA board behavior. When the CPU failed to boot (PC=0x00000000), this behavior was observed identically in both iverilog simulation and Vivado hardware. This correlation was instrumental in enabling rapid root cause analysis, as simulation iterations take seconds compared to the 10-minute FPGA build cycle.


The key factors that caused the simulation-to-hardware gap were not related to timing, clock domain crossing, or synthesis optimization—the usual suspects in FPGA debugging. Instead, they were related to file format compatibility, macro visibility, and transport delay behavior in simulation primitives. This finding underscores the importance of establishing a verified simulation environment before committing to hardware builds.


\subsection{ITCM/DTCM Initialization Complexity}

The discovery that both Vivado’s synthesis-time $readmemh and iverilog’s simulation-time $readmemh interpret whitespace-separated hex tokens as full memory words, combined with the E203’s use of different data widths for ITCM (64-bit) and DTCM (32-bit), created a subtle but critical initialization bug. The riscv64-unknown-elf-objcopy -O verilog command produces byte-level hex format, which is incompatible with word-level \$readmemh interpretation for memories wider than 8 bits.


The solution—a configurable-width hex format converter—highlights the importance of understanding the specific requirements of each tool in the build chain. This issue would not manifest in an ASIC flow where memory initialization is handled through different mechanisms.


\subsection{ILA Probe Design Strategy}

The initial ILA probe configuration monitored the ITCM-dedicated bus interface (ifu2itcm\_icb\_*), which was a natural choice given that the hello\_e203 program resides in ITCM. However, the CPU’s initial boot sequence begins in mask ROM (MROM), which is accessed through the BIU and system memory bus—not the ITCM interface. This mismatch caused the ILA to show no activity despite the CPU being active, leading to incorrect initial diagnosis.


The lesson for future ILA-based debugging is to include probes on all major bus interfaces during initial bring-up, rather than optimizing for the expected operational path.


\section{NICE CNN Accelerator Validation Status}\label{sec:5_2}

The NICE CNN accelerator test program successfully demonstrated:


The CPU correctly decodes and dispatches custom-0 opcode (0x0B) instructions to the NICE coprocessor.


The NICE interface correctly handles the four-channel handshake protocol (request, response, memory, CSR).


The PE array receives weight and activation data through the WLOAD and DLOAD instructions and responds to the COMP trigger.


However, the RSTAT result comparison did not match the expected software value. This discrepancy has been characterized and bounded—it affects only the final result readback, not the core instruction execution path. The ILA evidence confirms that all instructions complete without bus errors, pipeline stalls, or exception traps.


Several hypotheses for the RSTAT discrepancy are under investigation: (1) the WLOAD/DLOAD bank selection encoding (rs2[1:0]) may differ from the PE array’s internal addressing; (2) the INT8 data packing order within 32-bit registers may not match the byte lane mapping expected by the PE array; and (3) the accumulator readout path may require additional pipeline synchronization.


\section{Comparison with Project Objectives}\label{sec:5_3}

\begin{table}[htbp]
\centering
\caption{Project objective status summary}
\label{tab:5_1}
\begin{tabular}{lcc}
\hline
Objective & Status & Evidence \\ \hline\hline
Design 4×4 PE array with NICE interface & Achieved & RTL simulation RSTAT=19 \\ \hline
Integrate with E203 SoC & Achieved & Full-SoC simulation + FPGA build \\ \hline
FPGA bring-up pipeline & Achieved & 6 build modes, all timing-clean \\ \hline
hello\_e203 board validation & Achieved & UART output + ILA evidence \\ \hline
NICE accelerator board testing & Partially achieved & Instructions execute; RSTAT mismatch under investigation \\ \hline
Performance comparison (10x target) & Not yet measured & Requires working NICE readback \\ \hline
Accuracy validation (MNIST/LeNet-5) & Not yet performed & Requires working CNN application \\ \hline
\end{tabular}
\end{table}

\section{Engineering Lessons Learned}\label{sec:5_4}

Build verification infrastructure first. The ILA probe system, simulation environment, and hex format converter were essential enablers for all subsequent debugging. Without them, the four root causes would have been extremely difficult to isolate.


Simulation is not optional for complex SoC bring-up. Even though the final target is FPGA hardware, the ability to reproduce board behavior in simulation was the single most important factor in resolving the boot failure.


Document every configuration change. The interplay between E203\_FORCE\_BOOTROM\_BOOT, the prologue.tcl global defines, the fpga\_mem\_init.vh header, and the ITCM/DTCM hex file formats involved changes across multiple files in two repositories. Systematic documentation prevented confusion during iterative debugging.


Peripheral initialization order matters. The GPIO IOF configuration must precede UART initialization, and the UART baud rate divisor must be set before enabling the transmitter. These dependencies, while obvious in hindsight, caused several failed iterations when developing the NICE test program.


\section{Limitations}\label{sec:5_5}

The current work has several acknowledged limitations:


The CNN accelerator has been tested only with a 4×4 identity weight matrix and all-ones activation input. Full functional validation with realistic CNN weight data has not been performed.


Software-visible performance metrics (cycle count, speedup relative to CPU-only execution) have not been measured due to the incomplete RSTAT readback.


The MNIST/LeNet-5 accuracy validation, which was part of the original project scope, has not been completed.


The UART output issue in the NICE test program, while bypassed through LED-based indication, represents an unresolved integration issue that merits further investigation.


These limitations are clearly scoped and do not undermine the core contributions of the project: the complete FPGA bring-up pipeline, the systematic debugging methodology, and the demonstrated execution of NICE custom instructions on FPGA hardware.

